\section{Inversion and line-arrangement input}

Fix $p\in S$ and invert the remaining $N=n-1$ points about $p$.  A proper circle through $p$ becomes a nonradial line, while a line through $p$ becomes a radial line.  This transforms local circle incidences into line arrangements and permits the use of Melchior's inequality, ordinary-line bounds and Langer's incidence inequality \cite{kellymoser1958,csimasawyer1993,langer2003,dezeeuw2018}.

Let $\ell_j(p)$ denote the number of spanned lines in the inverted set containing exactly $j$ inverted points and put
\[
 W_p=\sum_{j\ge2}\frac{\ell_j(p)}{j+1}.
\]
If radial lines are retained, summing $W_p$ over $p\in S$ counts every generalized-circle block exactly once:
\begin{equation}\label{eq:weightedall}
 \sum_{p\in S}W_p=B.
\end{equation}
If radial terms are removed, the same weighting counts proper-circle blocks exactly once.

For $N$ real planar points with at most $R\le2N/3$ collinear, the pair identity, Melchior's inequality and Langer's inequality imply the exact dual estimate
\begin{equation}\label{eq:weightedLM}
 W\ge
 \frac{N^2R+2N^2+3NR-6N+18R+36}{36(R+1)}.
\end{equation}
Indeed, with
\[
 \alpha=\frac{R+2}{18(R+1)},\quad
 \beta=\frac{R+2}{6(R+1)},\quad
 \gamma=\frac{R-1}{36(R+1)},
\]
the coefficientwise slack for a $j$-point line factors as
\[
 \frac{(R-j)(j-3)(j-2)}{12(R+1)(j+1)}\ge0.
\]
This is the weighted Langer--Melchior certificate used in the global range and in the $n=16$ boundary case.  Simplicial-arrangement classifications and pseudoline inequalities are used only in explicitly identified finite branches \cite{cuntz2012,shnurnikov2016}.

